(1) 线性热导模型给由热端向冷端的热流率 \(\dot Q=G(T_h-T_c)>0\)，其中热导 \(G\) 单位为 \(\mathrm{W/K}\)。两热库温度由外部控制保持稳定，接触通道中的能流随温差反向而反向。 这里必须使用题设的平衡条件与参数范围，不应将该结果不加验证地推广到非平衡、非理想或高阶修正情形。
(2) 热端失热，熵变化率 \(-G(T_h-T_c)/T_h\)；冷端得热，熵变化率 \(+G(T_h-T_c)/T_c\)。两者相加得到 \(\dot S_{\rm gen}=G(T_h-T_c)^2/(T_hT_c)\ge0\)，温差不为零时严格大于零，满足热力学第二定律。
(3) 可逆热机在 \(T_h\) 吸热 \(\dot Q_h\)，向冷端排热 \(\dot Q_c=\dot Q_hT_c/T_h\)，故最大机械功率为 \(\dot W_{\max}=\dot Q_h(1-T_c/T_h)\)。直接导热没有产生机械功，因此温差对应的热力学可用能被不可逆传热所消耗。
(4) 设 \(T_c=T,T_h=T+\Delta T\)，\(\Delta T/T\ll1\)，则 \(\dot S_{\rm gen}=G(\Delta T)^2/[T(T+\Delta T)]=G(\Delta T)^2/T^2+O((\Delta T)^3)\)，无一阶温差项。热平衡 \(\Delta T\to0\) 时熵产生率趋零。
